\documentclass[10pt]{article}

\usepackage[utf8]{inputenc}

\usepackage[T1]{fontenc}

\usepackage{amsmath}

\usepackage{amsfonts}

\usepackage{amssymb}

\usepackage[version=4]{mhchem}

\usepackage{stmaryrd}

\usepackage{bbold}



\begin{document}

случае $p=$ char $k>0$ и $R$-порядок в алгебре кватернионов над $\mathbb{Q}$, разветвленной только в $p$ и $\infty$. Такие Э. к. существуют для всех $p$ и наз. с упе р сингул я р н ы м и; несуперсингулярные Э. к. в характеристике $p$ наз. обы кн о в ен ным и Ә. к.



Группа $X_{n}=$ Ker $n_{X}$ точек Э. к. $X$, порядок к-рых делит $n$, имеет следующую струкгуру: $X_{n} \approx(\mathbb{Z} / n \mathbb{Z})^{2}$, если $(n, \operatorname{char} k)=1$. При $p=\operatorname{char} k>0$ для обыкновенных Э.к. $X_{p^{m}} \approx \mathbb{Z} / p^{m} \mathbb{Z}$, а для суперсингулярных Ә.к. $X_{p} m \approx\{0\}$. Для простого $l \neq \operatorname{char} k$ Теüта модуль $T_{l}(X)$ изоморфен $\mathbb{Z}_{l}^{2}$.



Э. к. на д н з а мкну ты м и по ля ми. Пусть $X$ - Э. к. над произвольным полем $k$. Если множество $k$-рациональных точек $X(k)$ кривой $X$ непусто, то $X$ бирегулярно ивоморфна плоской кубич. кривой (1) с $a, b \in k($ char $k \neq 2,3)$. Бесконечно удаленная точка $P_{0}$ кривой (1) определена над $k$. Как и выше, можно определить групповую структуру на кривой (1), превращающую $X$ в одномерное абелево многообразие над $k$, а множество $X(k)$ в коммутативную группу с нулем $P_{0}$. Если $k$ конечно порождено над своим простым подполем, то $X(k)$ - группа с конечным числом образующих (т е о р е а М о р д е л л а - В е й л я).



Для любой Э. к. $X$ определено Якоби многообразие $J(X)$, являющееся одномерным абелевым многообразием над $k$. Э. к. $X$ является главным однородным пространством над $J(X)$. Если множество $X(k)$ непусто, то выбор точки $P_{0} \in X(k)$ задает изоморфизм $X \simeq J(X)$, при к-ром точка $P_{0}$ переходит в нуль групшы $J(X)$. В общем случае Э. к. $X$ и $J(X)$ изоморфны над конечным расширением поля $k$ (см. [1], [4], [13]).



Ә. к. над по ле м к о плексных чисе л. Э. к. $X$ нал $\mathbb{C}$ является комшактной римановой поверхностыю рода 1 и обратно. Групповая структура превращает $X$ в комплексную группу Ли, являющуюся одномерным комплексным тором $\mathbb{C} / \Lambda$, где $\Lambda$-решетка в комплексной плоскости $\mathbb{C}$. Обратно, любой одномерный комплексный тор является Ә. к. (см. [3]). С топологич. точки зрения Э. к.- двумерный тор.



Теория Э. к. над полем $\mathbb{C}$, по существу, эквивалентна теории әллиптич. функций. Отождествление тора $\mathbb{C} / \Lambda$ с Э. к. осуществляется следующим образом. Эллиптич. функции с данной решеткой периодов $\Lambda$ образуют поле, порожденное 8 -функцией Вейерштрасса (см. Вейерштрасса әллиптические бункции) и ее производной $8^{\prime}(z)$, к-рые связаны соотношением



\[

8^{\prime 2}=48^{3}-g_{2} 8^{\circ}-g_{3} .

\]



Отображение $\mathbb{C} \longrightarrow \mathrm{P}^{2}\left(Z \longrightarrow\left(1: 8(Z): 8^{\prime}(Z)\right)\right)$ индуцирует изоморфизм тора $\mathbb{C} / \Lambda$ и Э.к. $X \subset \mathrm{P}^{2}$ с уравнением $y^{2}=4 x^{3}-g_{2} x-g_{3}$. Отождествление Э. к. $X$, заданной уравнением (1), с тором $\mathbb{C} / \Lambda$ осуществляется с помощью криволинейных интегралов от голоморфной формы $\omega=$ $=\frac{d x}{y}$ и приводит к совнадению Э. к. $X$ с ее многообразием Якоби $J(X)$.



Описание множества всех Э. к. как торов $\mathbb{C} / \Lambda$ приводит к модулярной функциии $J(\tau)$. Две решетки $\Lambda$ и $\Lambda^{\prime}$ определяют изоморфные торы тогда и только тогда, когда они подобны, т. е. одна получается из другой умножением на комплексное число. Поэтому можно считать, что решетка $\Lambda$ порождена числами 1 и $\tau$ из $H=\{\tau \in \mathbb{C} \mid \operatorname{Im} \tau>0\}$. Две решетки с базисами $1, \tau$ и 1 , $\tau^{\prime}$ подобны тогда и только тогда, когда $\tau^{\prime}=\gamma(\tau)$ для нек-рого әлемента $\gamma$ модулярной группь Г. Модулярная функция



\[

J(\tau)=\frac{g_{2}^{3}}{g_{2}^{3}-27 g_{3}^{2}}

\]



наз. также а б с о л ю т ным ин в а р п а н т о м; $J(\tau)=J\left(\tau^{\prime}\right)$ тогда и только тогда, когда $\tau^{\prime}=\gamma(\tau)$ для нек-рого $\gamma \in \Gamma$ и функция $J: H / \Gamma \rightarrow \mathbb{C}$ осуществляет взаимно однозначное соответствие между классами изоморфных Э. к. над $\mathbb{C}$ и комплексными числами. Если $X=\mathbb{C} / \Lambda$, то $j(X)=1728 J(\tau)$.



Э. к. $X$ есть Э. к. с комплексным умножением тогда и только тогда, когда $\tau$ - мнимая квадратическая иррациональность. В этом случае $R$ - подкольцо конечного индекса в кольце целых алгебраич. чисел мнимого квадратичного поля $\mathbb{Q}(\tau)$. Ә. к. с комплексным умножением тесно связаны с полей классов теорией для мнимых квадратичных полей (см. [4], [8]).



А р ифме тика Э. к. Пусть $X$ - Э. к. над конечным полем $k$ из $q$ элементов. Множество $X(k)$ всегда непусто и конечно. Тем самым $X$ снабжается структурой одномерного абелева многообразия над $k$, a $X(k)-$ структурой конечной коммутативной группы. Порядок $A$ группы $X(k)$ удовлетворяет неравенству $\mid(q+1)-$ $-A \mid \leqslant 2 \sqrt{q}$. Многочлен $t^{2}-(q+1-A) t+q$ есть характеристич. многочлен $\Phi$ робениуса эндожорфизма, действующего на модуле Тейта $T_{l}(X), l \neq \operatorname{char} k$. Его корни $\alpha, \bar{\alpha}-$ комплексно сопряженные целые алгебраич. числа, по модулю равные $\sqrt{q}$. Для любого конечного расширения $k_{n}$ поля $k$ степени $n$ порядок групшы $X\left(k_{n}\right)$ ра̇вен $q^{n}+1-\left(\alpha^{n}+\overline{\alpha^{n}}\right)$. Дзета-функция Э. к. $X$ равна



\[

\left(1-q^{-s}\right)\left(1-q^{1-s}\right) /\left[1-(q+1-A) q^{-s}+q^{1-2 s}\right] .

\]



Для любого целого алгебраического $\alpha$, лежащего в нек-ром мнимом квадратичном поле (или в $\mathbb{Q}$ ) и по модулю равного $\sqrt{q}$, найдется такая Э. к. $X$ над $k$, что порядок групшы $X(k)$ равен $q+1-(\alpha+\bar{\alpha})$.



Пусть $k$ - поле $p$-адических чисел $\mathbb{Q}_{p}$ или его конечное алгебраич. расширение, $B$ - кольцо целых поля $k$, $X$ - Э. к. над $k$ и пусть множество $X(k)$ непусто. Групповая структура поревращает $X(k)$ в коммутативную компактную одномерную $Л u$ p-адическую гр уnny. Группа $X(k)$ двойственна по Понтрягину к Вейля - III amле ә pynne $W C(k, X)$. Если $j(X) \notin B$, то $X$ - кривая Тейта (cм. [1], [5]) и существует канонич. униформизация группы $X(k)$, аналогичная случаю поля $\mathbb{C}$.



Пусть $X-\ni . к$. над $\mathbb{Q}$ и множество $X(\mathbb{Q})$ непусто. Тогда $X$ бирегулярно изоморфна кривой (1) с $a, b \in \mathbb{Z}$. Из всех кривых вида (1) с пелыми $a$ и $b$, изоморфных $X$, выбирается такая, для к-рой абсолютная величина дискриминанта $\Delta$ минимальна. Кондуктор $N$ и $L$ функция $L(X, s)$ Э. к. $X$ определяются как формальные произведения локальных множителей



\[

\begin{gathered}

N=\Pi f_{p}, \\

L(X, s)=\Pi L_{p}(X, s)

\end{gathered}

\]



по всем простым $p$ (см. [1], [5], [13]). Здесь $f_{p}$ - нек-рая степень $p, L_{p}(X, s)$ - мероморфная функция комплексного переменного $s$, не имеющая ни нуля, ни полюса при $s=1$. Чтобы определить локальные множители, рассматривается редукция кривой $X$ по модулю $p(\neq 2,3)$ - плоская проективная кривая $X_{p}$ над полем вычетов $\mathbb{Z} /(p)$, заданная в афффинной системе координат уравнением



\[

y^{2}=x^{3}+\overline{a x}+\bar{b}(\bar{a}=a \bmod p, \bar{b}=b \bmod p) .

\]



Пусть $A_{p}$-число $\mathbb{Z} /(p)$-точек на $X_{p}$. Если $p$ не делит $\Delta$, то $X_{p}-\ni$. к. над $\mathbb{Z} /(p)$ п полагают



\[

f_{p}=: 1, \quad L_{p}(X, s)=1 /\left[1-\left(p+1-A_{p}\right) p^{-s}+p^{1-2 s}\right] .

\]



Если $p$ делит $\Delta$, то многочлен $x^{3}+\bar{a} x+\bar{b}$ имеет кратный корень и полагают



\[

L_{p}(X, s)=1 /\left[1-\left(p+1-A_{p}\right) p^{-s}\right], f_{p}=p^{2} \text { илй } p

\]



(в зависимости от того, является этот корень трехкратным или нет). Произведение (2) сходится в правой полу-





\end{document}